\section{Problem Statement}
\label{section:problem}
Our objective is to manipulate image contrast to enhance the sharpness and make image features more visually distinguishable. While image enhancement is inherently subjective, we have employed a fitness function from image enhancement literature to quantify this enhancement \cite{bhandari2020cuckoo,draa2014artificial,dos2009differential,braik2007particle,munteanu2004gray}. This fitness function evaluates the quality of an enhanced image. We use a transformation function to enhance the image, which is optimized using a metaheuristic approach. Thus, the input to our process is an image, and the output is an enhanced version of that image.

For simplicity, we present a formal definition of the problem in the context of a {\em grayscale} image. Let $I = f(x, y)$ represent a grayscale image, where $x$ and $y$ denote the pixel positions. The image $I$ has dimensions $M \times N$, thus $0 \leq x < M$ and $0 \leq y < N$. Assume there is a fuzzy logic-based transformation function $\transform{I}$ that modifies the gray value of each pixel, producing an enhanced image $I_e$. More formally,


% We give a more formal definition of the problem below in the context of a gray-scale image (for simplicity). Suppose a gray-scale image, $I = f(x,y)$, where $x$ and $y$ denote the pixel's positions of the image. Let the image $I$ have a dimension of $M \times N$. So,  $0\leq x < M$ and $0\leq y < N$. Now assume that there is a fuzzy logic-based transformation function $\transform{I}$ that transforms the gray value of each pixel of the image and outputs another image or enhanced image $I_e$. More formally,
\begin{equation} \label{eqn:transformation}
I_e = \transform{I} = \transform{f(x,y)}
\end{equation}
Thus, our task is to identify a suitable transformation function that visually enhances the resulting image $I_e$.









